\chapter{Introduction} \label{cha:introduction}
This chapter introduces the motivation behind this thesis, clearly defines the underlying problem and research question and presents an overview of the research.


\section{Motivation}
\ac{fl} represents a significant shift in how machine learning can be conducted, focusing on privacy and data security. Unlike traditional machine learning methods that often require the aggregation of data in a central location, \ac{fl} allows multiple participants to collaboratively train a model while keeping their data localized \cite{yang_federated_2019} \cite{li_survey_2023} \cite{kairouz_advances_2021}. 
\begin{figure}[htbp]
    \centerline{\includesvg[width=\columnwidth]{assets/fl.svg}}
    \caption[Traditional Federated Learning Process]{Traditional Federated Learning Process (based on \cite{jeromemetronome_federated_2019})}
    \label{fig:fl_process}
\end{figure}

\noindent The steps of the traditional \ac{fl} process are shown in Figure \ref{fig:fl_process}. The central server initializes a global model (Step 1) and distributes it to participating clients (Step 2). Each client trains the model locally using its own data (Step 3) and sends the updated model parameters back to the server. The server aggregates these updates, often by averaging, to improve the global model (Step 4). This cycle repeats for multiple rounds until the model converges \cite{zhang_survey_2021}. This is particularly beneficial in scenarios where data is decentralized, sensitive or subject to regulatory constraints like healthcare, finance or \ac{iot} \cite{li_survey_2023} \cite{zhang_survey_2021} \cite{kerkouche_privacy-preserving_2021}.

\ac{fl} has been extensively studied, but several challenges remain, particularly within edge computing environments, where devices may have varying computational capabilities and network conditions. Firstly, the traditional centralized aggregation of \ac{fl} raises concerns in terms of single points of failure and scalability issues in scenarios with limited bandwidth between nodes and a central server \cite{martinez_beltran_decentralized_2023} \cite{ko_asynchronous_2023}. Secondly, a significant concern involves the straggler effect, where resource disparities among the clients cause waiting times and delay the training process in synchronous systems resulting in reduced performance \cite{ko_asynchronous_2023} \cite{vu_straggler_2021}.

To address the single-point failure limitations of centralized FL, researchers have proposed \ac{dfl}, often utilizing block\-chains's inherent decentralization to enable edge devices to collaboratively build global models via \ac{p2p} communication \cite{martinez_beltran_decentralized_2023} \cite{nguyen_federated_2021}. To mitigate the problem resulting from the straggler effect, \ac{afl} has been introduced, facilitating the aggregation of local models without delays caused by slower-performing devices \cite{vu_straggler_2021} \cite{zang_general_2023} \cite{lu_differentially_2020}. 

Combining the benefits of asynchronous and decentralized approaches, many recent studies propose the use of \ac{dag} to store manage updates \cite{xue_energy_2022} \cite{lu_blockchain_2020} \cite{yuan_chainsfl_2021} \cite{cao_dag-fl_2021} \cite{xiao_db-fl_2024} \cite{liu_decentralized_2024} \cite{pvn_edge-enabled_2024} \cite{beilharz_implicit_2021} \cite{huang_personalized_2024} \cite{schmid_tangle_2020} \cite{cao_toward_2023}. \ac{dag}-based architectures provide specific advantages over traditional linear blockchain frameworks, particularly in edge computing scenarios. The intrinsic asynchronous properties of \ac{dag} networks effectively accommodate delayed transactions, aligning naturally with the operational requirements of asynchronous \ac{fl} \cite{ko_asynchronous_2023} \cite{zhao_evaluating_2019}. Each model update is a node in the \ac{dag}, and new updates reference previous ones, called "tips." The method for choosing these tips is the tip selection algorithm \cite{xiao_fast_2022}. However, the lack of a central authority introduces identity management challenges, making \ac{dag}-based \ac{fl} vulnerable to adversarial manipulation, particularly sybil attacks, that can degrade performance or introduce backdoors in the global model\cite{ko_asynchronous_2023} \cite{fung_limitations_2020}.

\section{Problem Statement}
Sybil attacks occur when an adversary creates multiple sybil identities (fake clients) to inject poisoned updates into a \ac{fl} system, aiming to inject poisoned gradients to degrade or backdoor the global model. These attacks distort learning outcomes and compromise model reliability \cite{fung_limitations_2020} \cite{jiang_sybil_2021}. Unlike traditional \ac{fl}, where extensive studies have been conducted to mitigate these kind of network attacks by detection or aggregation techniques through the central aggregation server \cite{jiang_sybil_2021} \cite{ghafourian_safl_2024} \cite{rodriguez-barroso_dynamic_2022} \cite{malecki_simeon_2021} \cite{garcia-marquez_krum_2025}, \ac{dag}-based \ac{fl} lacks a robust mechanism to verify participant legitimacy, making it highly susceptible to such attacks \cite{ko_asynchronous_2023}. Defenses specific to decentralized networks, such as \ac{pow} and \ac{pos}, are inefficient in edge environments due to high computational costs and potential collusion \cite{puthal_proof_2019}.

Although many \ac{dag}-based \ac{fl} studies discuss the tip selection algorithm as a natural effective solution to exclude underperforming or even poisoned models \cite{cao_dag-fl_2021} \cite{cao_toward_2023} \cite{yuan_chainsfl_2021}, none of them discusses this in sybil settings. Furthermore, tip selection as discussed in the existing studies might be able to identify poisoned model updates that aim to degrade model performance through accuracy measurements, but poisoned model updates that aim to backdoor the global model are not identifiable by accuracy measurements and are further amplified in a sybil setting \cite{andreina_baffle_2021} \cite{rodriguez-barroso_dynamic_2022}. This lack of consideration of network attacks in existing \ac{dag}-based \ac{fl} studies is additionally highlighted by the survey study from Seoyoung Ko et al. about \ac{dag}-based \ac{fl} which emphasizes the criticality to realize a secure solution that can defend against poisoning attacks (both model degradation and backdoors) and especially network attacks like sybil attacks \cite{ko_asynchronous_2023}.

Given the increasing adoption of \ac{fl} in edge environments like \ac{iot}, \ac{iiot}, \ac{iov} and \ac{iomt} \cite{zhang_survey_2021} \cite{li_survey_2023}
and the rising interest in \ac{dag}-based \ac{fl} to address decentralization and asynchronicity \cite{ko_asynchronous_2023}, ensuring reliable model aggregation even under adversarial settings is essential. Consider an industrial \ac{iot} network where edge devices in a smart factory collaboratively train an anomaly detection model using \ac{dag}-based \ac{fl}. Each device updates the model based on local sensor data and submits transactions to the \ac{dag} network. A malicious entity infiltrates the system, creating hundreds of sybil nodes that submit manipulated updates. These fake identities approve only each other’s transactions, ensuring that the compromised models dominate tip selection. As a result, the aggregated global model is corrupted, failing to detect critical machine failures and causing unexpected production downtime or safety hazards. In this scenario, the lack of sybil-resistance leads to operational failures and financial losses, demonstrating the urgent need for effective sybil attack prevention mechanisms in \ac{dag}-based \ac{fl}.

\section{Research Question}
Fung et al. state the following in their influential paper about the limitations of traditional federated learning in sybil settings:
"When each client’s training data is non-IID and has a unique distribution, we assume that honest clients can be distinguished from act-alike sybils by the diversity of their gradient updates" \cite{fung_limitations_2020}. The prerequisite of \ac{non-iid} data, which means that the statistical properties of the data points are not uniform across the dataset, holds especially in the domain of edge computing environments considered in this thesis, because of the heterogeneity among edge devices \cite{lu_federated_2024}. Based on the aforementioned statement, this thesis aims to answer the following research question:

\noindent\hrulefill\\
\noindent\textit{How can the distinguishability of honest clients and act-alike sybils by the diversity of their gradient updates be utilized during the tip selection of \ac{dag}-based \ac{fl} to reduce the influence of sybil poisoning attacks?}

\noindent\hrulefill\\
This research question directly addresses the significant gap highlighted in current \ac{dag}-based federated learning studies, specifically the absence of robust defenses against sybil attacks. By designing and evaluating a novel approach for \ac{dag}-based \ac{fl} that integrates gradient update similarity measurements in the tip selection algorithm, this thesis aims to measure the effectiveness of mitigating adversarial influence in sybil settings compared to traditional tip selection in \ac{dag}-based architectures.

\pagebreak
\section{Delimitations}
This thesis focuses on mitigating sybil poisoning attacks in \ac{dag}-based \ac{fl} by leveraging the distinguishability of honest clients and sybil nodes through gradient diversity in tip selection. While this research aims to enhance security in decentralized asynchronous federated learning, several constraints define its scope.

First, the study is limited to sybil poisoning attacks that manipulate model updates to degrade the model's performance or introduce backdoors, rather than other forms of adversarial behavior. Other adversarial tactics, including model inversion or membership inference are outside the scope of this work. 

Second, this thesis assumes an edge computing environment where devices exhibit computational heterogeneity and \ac{non-iid} data distributions. While these characteristics align with real-world \ac{iiot} and edge applications, findings may not be generalizable to other network settings. 

Third, the research focuses on tip selection mechanisms in \ac{dag}-based \ac{fl}. Alternative security solutions, such as cryptographic identity verification or consensus mechanisms like \ac{pow} or \ac{pos}, are not explored due to their computational overhead and inefficiencies in resource-constrained edge environments. 

Fourth, this study evaluates the proposed solution using simulations and controlled experiments rather than real-world deployments. While the simulation parameters are designed to reflect practical edge computing conditions, real-world variations in network latency, client availability, and attack sophistication may introduce additional challenges not captured in this analysis. 

Fifth, the thesis does not address privacy-preserving techniques such as differential privacy or secure multi-party computation. While privacy is an essential consideration in federated learning, the primary focus here is on sybil resilience rather than data confidentiality. 

Finally, this research does not aim to provide a universally optimal tip selection algorithm but rather investigates the effectiveness of gradient diversity as a distinguishing factor for sybil mitigation. The approach will be compared against baseline tip selection strategies, but broader comparisons with all existing tip selection algorithms fall beyond the scope of this work.

\pagebreak
\section{Thesis Structure}
The thesis is structured as follows:

\paragraph{Chapter \ref{cha:introduction}} contains the introduction describing the motivation, research problem and the research question guiding this thesis. 

\paragraph{Chapter \ref{cha:extended_background}} dives deeper into the foundations of \ac{fl}, its challenges, \ac{dfl}, \ac{afl}, \ac{dag}-based \ac{fl} and sybil attacks. Furthermore existing research is reviewed, gaps identified and this thesis positioned in the landscape of current research. 

\paragraph{Chapter \ref{cha:methodology}} explains the selected approach of \ac{dsr} and the used research strategies and methods. Furthermore, the method choice in relation to the research methods applied in related literature is discussed and ethical considerations disclosed. 

\paragraph{Chapter \ref{cha:sydag-fl}} concretizes the problem addressed though a threat model of sybil attacks on \ac{dag}-based \ac{fl}. Based on this and the literature review, the requirements for the approach are stated. The approach is then formally presented.

\paragraph{Chapter \ref{cha:sydag-fl_evaluation}} evaluates the created approach. First the conducted experiment is described in detail an secondly the results of the experiment are presented.

\paragraph{Chapter \ref{cha:conclusion}} concludes the thesis by firstly interpreting and discussing the results, secondly summarizing and outlining the practical and theoretical significance of the research contributions, thirdly discussing the ethical and social aspects of the use of the approach and, fourthly suggesting directions for future work.